\subsection{关于相对可分代数闭包的补充}

定理 1（Zariski）。--- 设 K 是一个域，L 是 K 的一个扩域， \(K_{1}\) 是 K 在 L 中的相对可分代数闭包（V，第 44 页）， \(K_{2}\) 是 K 在 L 中的相对代数闭包（V，第 19 页），以及 \((\mathbf{X}_{i})_{i\in I}\) 是一组未定元。那么 \(\mathrm{K}_{1}(\mathbf{X}_{i})_{i\in I}\) 是 \(\mathrm{K}(\mathbf{X}_{i})_{i\in I}\) 在 \(\mathrm{L}(\mathbf{X}_{i})_{i\in I}\) 中的相对可分代数闭包，且 \(\mathrm{K}_{2}(\mathbf{X}_{i})_{i\in I}\) 是 \(\mathrm{K}(\mathbf{X}_{i})_{i\in I}\) 在 \(\mathrm{L}(\mathbf{X}_{i})_{i\in I}\) 中的相对代数闭包。

\begin{enumerate}
\def\labelenumi{\Alph{enumi})}
\tightlist
\item
  设 E 是一个域，F 是 E 的一个扩域，并且 F 中每个在 E 上可分代数的元素都属于 E。令 u 为 \(F(X)\) 中的一个在 \(E(X)\) 上可分代数的元素；我们将证明 u 属于 \(E(X)\) 。在 \(F(X)\) 中存在两个互素的多项式 P 和 Q，使得 u = P/Q，并且我们可以取 Q 为首一多项式。设 S 为 F 的由 P 和 Q 的系数组成的有限子集， \(F_{0} = E(S)\) ，并令 \(\Delta\) 为 \(\mathbf{F}_0\) 到 \(\mathbf{F}_0\) 中的一个 E-导子。设 \(D\) 为 \(\mathrm{F}_0(\mathbf{X})\) 到自身的、与 \(\Delta\) 在 \(\mathbf{F}_0\) 上一致且把 \(X\) 映为 0 的导子（V，第 128 页，命题 3）。
\end{enumerate}

由于 \(u \in F_0(X)\) 在 \(E(X)\) 上是可分代数的，且 \(D\) 在 \(E(X)\) 上为零，我们有 \(D(u) = 0\) （V，第 129 页，命题 4），由此 \(D(P).Q = P.D(Q)\) 。由于 \(P\) 和 \(Q\) 是互素的，我们得出结论 \(Q\) 整除 \(D(Q)\) （IV，第 13 页，推论 4）。现在 \(Q\) 可以写成形式

\[
\mathrm{Q} (\mathrm{X}) = \mathrm{X} ^ {n} + a _ {1} \mathrm{X} ^ {n - 1} + \dots + a _ {n - 1} \mathrm{X} + a _ {n}
\]

其中 \(a_1, \ldots, a_n\) 属于 \(\mathbf{F}_0\) ；由于 \(\mathbf{D}(x) = 0\) ，因此我们有

\[
\mathrm{D} (\mathrm{Q}) = \Delta (a _ {1}) \mathrm{X} ^ {n - 1} + \dots + \Delta (a _ {n - 1}) \mathrm{X} + \Delta (a _ {n})\tag{1}
\]

由此 \(\deg D(Q) < \deg Q\) 。由于 \(Q\) 整除 \(D(Q)\) ，这只有在 \(D(Q) = 0\) 时才可能。但此时 \(D(P) = 0\) ，因为 \(D(P) \cdot Q = P \cdot D(Q)\) 。现在 (1) 以及对于 \(P\) 的一个类似公式表明 \(\Delta\) 零化系数集合 \(S\) （\(P\) 和 \(Q\) 的），从而 \(\Delta = 0\) ，因为 \(F_0 = E(S)\) 。根据 \(V\) ，第 135 页，推论 2，有限生成扩张 \(F_0\) （\(E\) 的）因此是可分代数的；根据关于 \(E\) 和 \(F\) 的假设我们有 \(F_0 = E\) ，所以最终 \(u \in E(X)\) 。

\begin{enumerate}
\def\labelenumi{\Alph{enumi})}
\setcounter{enumi}{1}
\item
  现在假设 E 在扩域 F 中是相对代数封闭的，并记 p 为 E 的特征指数。设 u 是 \(F(X)\) 中的一个在 \(E(X)\) 上代数的元素。存在一个整数 \(f \geqslant 0\) 使得 \(v = u^{p^f}\) 在 \(E(X)\) 上是可分代数的（V，第 44 页，命题 13）。由 A) 我们因此有 \(v \in E(X)\) 。存在 u 的形如 P/Q 的唯一表示，其中 P 和 Q 是 F{[}X{]} 中互素的多项式且 Q 为首一的；我们有类似的分解 \(v = P_1 / Q_1\) ，其中 \(P_1\) 和 \(Q_1\) 在 E{[}X{]} 中互素，且 \(Q_1\) 是首一的。由此可知 \(P_1 / Q_1 = P^{p^f} / Q^{p^f}\) ；多项式 \(P^{p^f}\) 和 \(Q^{p^f}\) 在 F{[}X{]} 中互素（IV，第 13 页，推论 6），如同 \(P_1\) 和 \(Q_1\) 一样，且 \(Q^{p^f}\) 是首一的。由此可得 \(P^{p^f} = P_1 \in E[X]\) 且 \(Q^{p^f} = Q_1 \in E[X]\) 。因此 P 和 Q 的系数在 E 上是 p-根式的，从而属于 E，因为 E 在 F 中是相对代数封闭的。我们因此有 \(P \in E[X]\) , \(Q \in E[X]\) 并且最终 \(u \in E(X)\) 。这证明了 \(E(X)\) 在 F(X) 中是相对代数封闭的。
\item
  我们使用定理 1 的记号。由于 \(\mathbf{K}_1\) 是 \(\mathbf{K}\) 的一个可分代数扩张，扩张 \(\mathbf{K}_1(\mathbf{X}_i)_{i\in I}\) （\(\mathbf{K}(\mathbf{X}_i)_{i\in I}\) 的扩张）是代数的且可分的（V，第 39 页，命题 6）。此外， \(L\) 的每个在 \(\mathbf{K}_1\) 上代数且可分的元素都属于 \(\mathbf{K}_1\) （V，第 44 页，命题 13, a)）。设 \(J\) 是 \(I\) 的一个有限子集；通过对 \(J\) 的基数的直接归纳，我们由 \(A\) ) 推出， \(\mathrm{L}(\mathbf{X}_i)_{i\in J}\) 的每个在 \(\mathbf{K}(\mathbf{X}_i)_{i\in J}\) 上代数且可分的元素都属于 \(\mathbf{K}_1(\mathbf{X}_i)_{i\in J}\) 。最后设 \(u\) 是 \(\mathrm{L}(\mathbf{X}_i)_{i\in I}\) 的一个在 \(\mathbf{K}(\mathbf{X}_i)_{i\in I}\) 上代数且可分的元素；存在一个有限子集 \(J\) （\(I\) 的），使得 \(u\) 属于 \(\mathrm{L}(\mathbf{X}_i)_{i\in J}\) 且在 \(\mathbf{K}(\mathbf{X}_i)_{i\in J}\) 上是代数且可分的；根据前述内容， \(u\) 属于 \(\mathbf{K}_1(\mathbf{X}_i)_{i\in J}\) ，更不用说属于 \(\mathbf{K}_1(\mathbf{X}_i)_{i\in I}\) 了。
\end{enumerate}

我们由此从 \(A\) 推出 \(\mathrm{K}_1(\mathbf{X}_i)_{i\in I}\) 是 \(\mathrm{K}(\mathbf{X}_i)_{i\in I}\) 在 \(\mathrm{L}(\mathbf{X}_i)_{i\in I}\) 中的相对可分闭包；同样地，由 \(B\) 推出 \(\mathrm{K}_2(\mathbf{X}_i)_{i\in I}\) 是 \(\mathrm{K}(\mathbf{X}_i)_{i\in I}\) 在 \(\mathrm{L}(\mathbf{X}_i)_{i\in I}\) 中的相对代数闭包。

